\documentclass[14pt,a4paper]{extarticle}
\usepackage[left=3cm,right=2cm,top=2cm,bottom=2cm,headheight=18pt]{geometry}
\usepackage{fontspec}
\setmainfont{Liberation Serif}
\setsansfont{Liberation Serif}
\setmonofont{Liberation Mono}
\usepackage{unicode-math}
\setmathfont{texgyretermes-math.otf}
\usepackage[english]{babel}
\usepackage{amsmath,graphicx,booktabs,longtable,array,enumitem}
\usepackage[normalem]{ulem}
\usepackage{titlesec,fancyhdr}
\usepackage[hidelinks,unicode]{hyperref}
\usepackage{url}
\urlstyle{same}
\setlength{\parindent}{1.27cm}
\setlength{\parskip}{0pt}
\linespread{1}
\setlength{\emergencystretch}{3em}
\clubpenalty=10000
\widowpenalty=10000
\titleformat{\section}{\normalfont\bfseries}{\thesection.}{0.5em}{}
\titlespacing*{\section}{0pt}{1.1em}{0.5em}
\setlist{nosep,leftmargin=1.27cm}
\pagestyle{fancy}
\fancyhf{}
\renewcommand{\headrulewidth}{0pt}
\renewcommand{\footrulewidth}{0pt}
\fancyhead[C]{\ifnum\value{page}>1\thepage\fi}
% Information-analytical material: O‘RQ-682 appended methodology §83 (11–14 pt), contents §84, comparison table §85.
\renewcommand{\normalsize}{\fontsize{14bp}{16.8bp}\selectfont}
\normalsize
\newcommand{\tsize}{\fontsize{12bp}{14.4bp}\selectfont}
\newcolumntype{L}[1]{>{\raggedright\arraybackslash}p{#1}}
\setlength{\LTcapwidth}{\linewidth}
\begin{document}
\begin{flushright}Material 2\end{flushright}
\begin{center}\textbf{EXPLANATORY NOTE}\par
\textbf{to the proposals for inclusion in the draft normative legal act being prepared to implement paragraph 14 of Decree of the President of the Republic of Uzbekistan No. PF-193 of 11 September 2026}\end{center}

\noindent Edition of 1 October 2026. English companion translation; the Russian edition is primary. The structure follows paragraph 84 of the Unified Methodology appended to the Law “On Normative Legal Acts”.

\section{Basis for preparation and necessity}
Subparagraph “a” of paragraph 13 of Decree PF-193 provides that from 1 January 2027 need is studied and families are included in the Social Registry through a multidimensional assessment system that takes into account, \textit{in addition to the family’s financial condition}, chronic illness, disability, external labour migration and large family size. Subparagraph “b” provides for a point-based scoring of socio-economic condition and assignment of families to categories by its results. Paragraph 14 instructs the National Agency for Social Protection under the President of the Republic of Uzbekistan (hereinafter, the Agency), jointly with the Ministry of Economy and Finance, to submit the draft normative legal act to the Presidential Administration by 1 December 2026\cite{pf193}.

The need for the proposed norms is supported by the following.

\textbf{First, the decision is made automatically.} Under stage 6 of the Scheme for Including Families in the Social Registry (Annex 1 to the Regulation on the Procedure for Maintaining the Social Registry), the “Ijtimoiy reyestr” information module recognizes or refuses a family and assigns its category “automatically”\cite{res35}. In a point-based system the category depends on which side of a threshold the score falls. Every threshold, its equality case and every exception therefore directly determine a family’s access to assistance.

\textbf{Second, the law requires an automated decision to be explained.} Article 24 of the Law “On Personal Data” obliges the owner and/or operator to explain to the data subject the procedure for a decision based solely on automated processing and its possible legal consequences, to give the opportunity to object, and to notify the result of considering the objection in writing within ten days\cite{pdlaw}. Point 37 of the Regulation provides only an SMS about inclusion or refusal\cite{res35}.

\textbf{Third, practice shows unjustified refusals.} The Government portal reported on 17 July 2026 that in 184 mahallas all 2,000 applications received in one month were refused without grounds; appeals to the President on Registry inclusion rose one and a half times since the start of the year, from 30,000 to 46,000; and the Prosecutor General’s Office was instructed to strictly oversee proper maintenance of the Registry\cite{govuz}. Such decisions can be detected and corrected only if the grounds of each decision are retained and reproducible.

\textbf{Fourth, four inconsistencies have been demonstrated in the current calculation rules.} Three of them (points 26, 27 and 33 of the Regulation) are removed by the amendments in Section II of Material 1 (Table 1); the fourth, concerning families in difficult life circumstances, requires the developer’s decision and is set out in Section~\ref{sec:q}. In addition, point 37 of the Regulation does not secure the requirements of Article 24 of the Law “On Personal Data”. Because the family’s financial condition remains a criterion from 1 January 2027, the income calculation rules continue to determine the assessment result until replaced.

\begingroup\tsize
\begin{longtable}{L{0.24\linewidth} L{0.43\linewidth} L{\dimexpr0.33\linewidth-6\tabcolsep}}
\caption*{Table 1. Demonstrated inconsistencies in the current Regulation}\\
\toprule
Rule & Demonstrated result & Proposed correction\\\midrule\endhead
Point 27: household plot area $TEM=UM-QM-\min(0.2\,UM;\,200)$ & Where buildings cover more than 80\,\% of a plot of up to 1,000 m², the area becomes negative. Example: $UM=100$, $QM=90$ gives $-10$ m² and $-4{,}760$ soum of monthly normative income (minimum wage 1,360,000 soum, mahalla coefficient 1). A fully built-up 1,000 m² plot gives $-200$ m² and $-95{,}200$ soum per month. Negative income understates the family’s income. & Section II, point 2: a negative result is taken as $TEM=0$.\\\midrule
Point 32, subpoint “a” of point 33, subpoints “a” and “b” of point 39 & Where average monthly income per member \textit{equals} minimum consumption expenditure (715,000 soum in 2026), the refusal ground “higher than expenditure” does not apply and, absent other refusal grounds, the family is recognized under point 32; but categories “a” and “b” of point 39 and point 3 of Decree PF-258 require income \textit{below} expenditure. A first-time applicant without a “mahalla seven” recommendation is assigned to no category. & Section II, point 3: refusal where income is \textit{equal to} or higher than expenditure, aligning the Regulation with Decree PF-258.\\\midrule
Point 26: normative income for remittances & If the transfer is compared as a monthly average, for countries listed in Annex 6 a transfer of $2{,}719{,}999$ soum is counted as $4{,}080{,}000$ (three minimum wages), while $2{,}720{,}000$ is counted as $2{,}720{,}000$. A one-soum increase in the actual transfer reduces counted income by $1{,}360{,}000$ soum. The comparison period (month or three months) and the comparison where several family members are abroad are not stated. & Section II, point 1: normative income applies where the monthly average transfer from a family member is below the normative income established for that member; counted income does not fall as the transfer rises. For transfers from Annex 6 countries between $2M$ and $3M$ per month, counted income rises to $3M$, as smaller transfers are already counted today.\\\midrule
Point 37 & Only an SMS about inclusion or refusal is provided; explanation of the automated decision procedure, its legal consequences and the procedure for protecting rights (Article 24, part three, of the Law “On Personal Data”) is not secured. & Section II, point 4: decision grounds, legal consequences, objection and rights-protection procedure through the personal account and Inson centres.\\
\bottomrule
\end{longtable}\endgroup

All figures in the table are reproduced by the program in the electronic supplement to Material 3 (39 automated checks pass; independent arithmetic control of 4,105 constructed cases of the area formula, of which 620 give a negative area).

\section{Author of the proposal}
The proposal was prepared by citizen Abduxoliq Akmal ugli Ashuraliyev, independent researcher (ORCID 0009-0003-5482-5526). The official developers of the draft under paragraph 14 of Decree PF-193 are the Agency jointly with the Ministry of Economy and Finance. Article 23 of the Law “On Normative Legal Acts” requires the developer to study results of scientific research and appeals of individuals, and to generalize and use proposals of citizens and recommendations of scientists and specialists\cite{orq682}.

\section{Purposes, tasks and main provisions}
The purpose is to ensure that each family’s assessment result follows only from approved rules, can be reproduced from retained data and is explained to the applicant. The tasks are: remove uncertainty at threshold values; check software before changes are applied; retain the grounds of each result; and eliminate the four demonstrated inconsistencies in the current Regulation.

\begingroup\tsize
\begin{longtable}{L{0.17\linewidth} L{0.40\linewidth} L{\dimexpr0.43\linewidth-6\tabcolsep}}
\toprule
Points of Material 1 & Content & Justification\\\midrule\endhead
Section I, point 1 & Terms & Methodology paragraph 8: repeatedly used terms are defined at the beginning.\\
Points 2–5 & Source, unit, period and scoring of each criterion; category at equality; order of exceptions; calculation examples & Without these elements the same income can be counted differently (Table 1, point 26). The current Regulation already contains a calculation example (point 30); Chile’s published methodology states thresholds as “equal to or greater than” and gives a precedence rule when several checks are triggered\cite{chile}.\\
Point 6 & Software applies only established conditions & A new refusal ground is established only by a normative act, not by a software setting.\\
Point 7 & Non-receipt of a reply is not absence of income or property & Separates a technical failure from a fact; current point 48 of the Regulation already allows a temporary algorithm change during failures\cite{res35}.\\
Points 8–11 & Comparison with control calculations before changes are applied; mandatory content of the example set; documentation within existing commissioning documents & The errors in Table 1 lie precisely at thresholds and in exceptions. Agreement of totals on random examples does not detect them.\\
Point 12 & Retention of methodology and software versions, values, score and threshold & Allows every decision affected by an error to be found and an objection to be answered on the merits.\\
Point 13 & Marking and review of results under a temporary algorithm & A temporary algorithm change does not become a hidden permanent rule.\\
Point 14 & Decision grounds and legal consequences provided to the applicant & Implements Article 24, part three, of the Law “On Personal Data”; the ten-day period for answering an objection is set by part four of that Article and is not repeated.\\
Point 15 & Special data of other family members only as the fact that a criterion applied & Health and conviction data are special personal data (Article 25 of the Law “On Personal Data”); other family members are separate data subjects.\\
Point 16 & Review of all affected decisions & An error is corrected for all families, not only for those who complained.\\
Point 17 & Financing & Section~\ref{sec:fin} of this note.\\
Section II, points 1–4 & Amendments to points 26, 27, 33 and 37 of the Regulation & Table 1 and Section~\ref{sec:cmp}.\\
\bottomrule
\end{longtable}\endgroup

\section{Regulation at the level of law}
The proposals relate to the subordinate act mandated by paragraph 14 of Decree PF-193. No draft law is proposed.

\section{Expected results}
\begin{enumerate}[leftmargin=1.27cm]
\item No household-plot area calculation gives a negative value (at present this occurs where $QM>0.8\,UM$ for plots up to 1,000 m² and where $QM>UM-200$ for plots over 1,000 m²).
\item For each threshold of the methodology the category at equality is determined; a family with income equal to minimum consumption expenditure receives a determinate result.
\item A family’s counted income does not fall when its actual remittance increases.
\item Each assessment result can be reproduced from retained data; when an error is found, all affected decisions are identified.
\item The applicant receives the grounds of the automated decision, information on its legal consequences and the objection and rights-protection procedure, as required by Article 24, part three, of the Law “On Personal Data”.
\end{enumerate}

\section{Foreign experience}
Chile’s Ministry of Social Development and Family publishes the procedure for calculating the socio-economic rating of the Social Household Registry, approved by Resolution No. 082 of 15 November 2023, with worked calculation examples, thresholds stated as “equal to or greater than” and an explicit precedence rule when several checks are triggered for a household\cite{chile}. The World Bank notes that outdated registry data lead to both inclusion and exclusion errors, and that grievance redress mechanisms enhance transparency and accountability\cite{wb2025}.

\section{Description of content and a question for the developer’s decision}\label{sec:q}
Section I of Material 1 (17 points, each containing one legal norm in accordance with Methodology paragraph 22) is included in the act being prepared. Section II (4 points) amends the current Regulation; if the new act replaces the Regulation, the same rules are included in its norms.

\textbf{Question for the developer’s decision.} For families in difficult life circumstances, points 34–35 of the Regulation allow inclusion in the “family at the poverty boundary” category where income exceeds minimum consumption expenditure $P$ and income reduced by 25\,\% does not exceed $1.5P$, that is, where $P<D\le2P$. Subpoint “v” of point 39 of the Regulation and subpoint “v” of point 3 of Decree PF-258 describe the same category by the interval from $P$ to $1.5P$. Under either reading of the interval, some families included under point 35 fall outside it:
\begin{itemize}[leftmargin=1.27cm]
\item if the interval applies to ordinary income, a family with income of $1{,}287{,}000$ soum ($1.8P$) is included under point 35 but falls outside the interval;
\item if the interval applies to reduced income, a family with income of $858{,}000$ soum ($1.2P$; reduced income $643{,}500<P$) is included under point 35 but falls outside the interval.
\end{itemize}
Resolving this requires choosing which income and which interval apply to this procedure, with coordinated amendment of the Regulation and, if necessary, Decree PF-258. Points 3–4 of Section I prevent such an inconsistency in the new methodology.

\section{Amendments to other acts}
Section II amends the Regulation on the Procedure for Maintaining the Social Registry (Annex 1 to Cabinet of Ministers Resolution No. 35 of 29 January 2026). The amendment of point 33 aligns the Regulation with point 3 of Decree PF-258 and does not require amending the Decree. Amendment of Decree PF-258 may be needed only for the question in Section~\ref{sec:q}.

\section{Justification of references}
Points 9, 15 and 16 of Section I refer to points 8, 14 and 12 of the same Section to avoid repeating their text (Methodology paragraph 13). Point 8 refers to point 13 so that the comparison does not delay a temporary algorithm change during data-exchange failures. Point 11 of Section I refers to the established procedure for commissioning and changing information systems without naming a specific act: references to specific acts of lower legal force are not permitted (Methodology paragraph 11), and the type of the act being prepared is determined by the developer. Section II names the amended act and its points (Methodology paragraph 69).

\section{Financial and economic justification}\label{sec:fin}
Implementation of the proposals \textbf{requires no additional expenditure from the State budget}:
\begin{enumerate}[leftmargin=1.27cm]
\item Section I uses the existing YAMIH system, the existing commissioning procedure and existing channels — the personal account on the Unified Portal of Interactive Public Services and Inson centres. No new bodies, posts, information systems or application stages are created.
\item The set of control examples and the program calculating them have already been prepared by the author and are provided free of charge with Material 3.
\item The amendments to points 26, 27 and 33 do not widen the circle of assistance recipients: a family’s counted income does not fall, and a case for which no category is established is removed. The amendment to point 26 raises to $3M$ the counted income of families with transfers from Annex 6 countries between $2M$ and $3M$ per month, as smaller transfers are already counted today; some such families may lose eligibility. If the developer considers this undesirable, the non-monotonicity can be removed differently — by setting a normative income of $2M$ for these countries; that option widens the circle of recipients and requires a cost estimate. The amendment to point 37 uses existing channels.
\end{enumerate}
The exact number of families affected by each amendment can be determined only by the Agency on its data using the attached program, without transferring personal data outside the Agency.

\section{Discussion and coordination}
The proposal is submitted under Article 20 of the Law “On Normative Legal Acts” and Article 3 of the Law “On Appeals of Individuals and Legal Entities”. Discussion on the portal and coordination with interested bodies are conducted by the developer. Article 23 of the Law “On Normative Legal Acts” requires materials of the discussion of a draft to be submitted to the adopting body together with the draft; part 2 of Article 21 provides for placing information on the preparation of the most important socio-economic drafts on the portal, and Article 24 provides for placing the draft on the portal before it is submitted to the adopting body\cite{orq682}. As of 1 October 2026, none of the 40 drafts placed on the discussion portal since 11 September 2026 is the draft under paragraph 14 of Decree PF-193\cite{portal}.

\section{Answers to possible objections}
\begingroup\tsize
\begin{longtable}{L{0.32\linewidth} L{\dimexpr0.68\linewidth-4\tabcolsep}}
\toprule
Objection & Answer\\\midrule\endhead
The current rules will be replaced from 2027 & Financial condition remains a criterion (Decree PF-193, paragraph 13, subparagraph “a”). Until replacement the Regulation applies; on replacement the same corrections are included in the new norms. Section I concerns the new methodology itself.\\\midrule
Software is already tested & Section I uses the existing commissioning procedure (point 11). It requires only a minimum content of the example set: each threshold, equality to it and each exception (point 9). If existing tests cover this, there is no additional work.\\\midrule
Existing powers suffice (Article 20 of the Law “On Normative Legal Acts”) & The proposals are made to an already mandated draft, not to initiate a new act. The inconsistencies in Table 1 are in the text of a normative act and can be removed only by amending that text. The right to an explanation of an automated decision is established by law and requires a defined channel.\\\midrule
The number of affected families is unknown & The inconsistencies follow from the text of the norms and exist for any data. The number of affected families determines the scale, not the existence of the error; the Agency can determine it with the attached program.\\\midrule
Costs & Section~\ref{sec:fin}: no additional budget expenditure is required.\\\midrule
Disclosure of information & Only information about the applicant’s own family is disclosed; special personal data of other family members are shown only as the fact that a criterion applied (Article 25 of the Law “On Personal Data”); source code and security information are not disclosed.\\
\bottomrule
\end{longtable}\endgroup

\section{Comparison table}\label{sec:cmp}
The Regulation is officially published in the state language; therefore the current and proposed wording is given in Uzbek, with a translation of the changed words in the justification column. Changed words of the current wording are italic and underlined; new words are bold (Methodology paragraph 85).

\begingroup\tsize
\begin{longtable}{L{0.33\linewidth} L{0.33\linewidth} L{\dimexpr0.34\linewidth-6\tabcolsep}}
\toprule
Current wording & Proposed wording & Justification\\\midrule\endhead
26-band, birinchi xatboshi: «…pul oʻtkazmalari miqdori aniqlanmasa yoki \textit{\uline{oʻtkazmalar miqdori mehnatga haq toʻlash eng kam miqdorining ikki baravari miqdoridan}} kam boʻlsa, u holda … quyidagicha normativ daromad qoʻllaniladi:» & «…pul oʻtkazmalari miqdori aniqlanmasa yoki \textbf{chet eldagi oila aʼzosidan kelgan oʻtkazmalarning bir oylik oʻrtacha miqdori ushbu bandda unga nisbatan belgilangan normativ daromaddan} kam boʻlsa, u holda … quyidagicha normativ daromad qoʻllaniladi:» & “less than twice the minimum wage” becomes “monthly average from a family member abroad less than the normative income established for that member”. Removes the fall in counted income when a transfer increases, the unspecified comparison period and the comparison where several members are abroad (Table 1).\\\midrule
27-band: «TEM = UM — QM — (UM × 20\%) … (UM × 20 \%) koʻrsatkichi 200 metr kvadratdan oshmagan holda hisob-kitob uchun qabul qilinadi.» & Oltinchi xatboshidan keyin: «\textbf{Formula boʻyicha aniqlangan TEM qiymati manfiy boʻlsa, TEM nolga teng deb qabul qilinadi.}» & New paragraph: “If the TEM value is negative, TEM is taken as zero.” Excludes negative area and negative normative income (Table 1).\\\midrule
33-band «a» kichik bandi: «…bir oylik jami oʻrtacha daromadi minimal isteʼmol xarajatlari miqdori\textit{\uline{dan}} yuqori boʻlganda.» & «…bir oylik jami oʻrtacha daromadi minimal isteʼmol xarajatlari miqdori\textbf{ga teng yoki undan} yuqori boʻlganda.» & “higher than expenditure” becomes “equal to or higher than expenditure”. A family with income equal to expenditure receives a determinate result consistent with PF-258 point 3 (Table 1).\\\midrule
37-band: «Oilaning Reyestrga kiritilganligi yoki Reyestrga kiritish rad etilganligi haqida ariza beruvchining mobil telefon raqamiga SMS-xabar yuboriladi.» & Ikkinchi xatboshi qoʻshiladi: «\textbf{Bir vaqtning oʻzida ariza beruvchiga Yagona interaktiv davlat xizmatlari portalidagi shaxsiy kabinet orqali, uning murojaatiga koʻra esa “Inson” ijtimoiy xizmatlar markazida qaror qabul qilinishiga asos boʻlgan maʼlumotlar, qoʻllanilgan normativ daromadlar, qaror qabul qilish tartibi, qarorning yuridik oqibatlari, unga eʼtiroz bildirish va oʻz huquqlarini himoya qilish tartibi toʻgʻrisidagi maʼlumotlar taqdim etiladi. Boshqa oila aʼzolarining shaxsga doir maxsus maʼlumotlari faqat tegishli mezon qoʻllanilganligi fakti sifatida koʻrsatiladi.}» & New paragraph: grounds of the decision, normative incomes applied, decision procedure, its legal consequences, objection and rights-protection procedure; special data of other family members only as the fact that a criterion applied. Implements Article 24, part three, and Article 25 of the Law on Personal Data.\\
\bottomrule
\end{longtable}\endgroup

\begingroup\tsize
\begin{thebibliography}{9}
\bibitem{pf193} Decree of the President of the Republic of Uzbekistan No. PF-193 of 11 September 2026, paragraphs 13–14. \url{https://lex.uz/docs/-8472526}.
\bibitem{res35} Resolution of the Cabinet of Ministers No. 35 of 29 January 2026, Annex 1 (Regulation on the Procedure for Maintaining the Social Registry), points 26–27, 32–35, 37, 39, 48 and Annex 1 to the Regulation. \url{https://lex.uz/docs/-8027513}.
\bibitem{pf258} Decree of the President of the Republic of Uzbekistan No. PF-258 of 26 December 2025, paragraph 3 and Annex 1. \url{https://lex.uz/docs/-7952180}.
\bibitem{pdlaw} Law of the Republic of Uzbekistan “On Personal Data” No. O‘RQ-547 of 2 July 2019, Articles 24 and 25. \url{https://lex.uz/docs/-4396419}.
\bibitem{orq682} Law of the Republic of Uzbekistan “On Normative Legal Acts” No. O‘RQ-682 of 20 April 2021, Articles 20–24; appended Unified Methodology, paragraphs 8, 11, 13, 22, 69, 84–85. \url{https://lex.uz/docs/-5378966}.
\bibitem{govuz} Government portal. “Mahalla odamlar uchun ishonch, tayanch va imkoniyatlar markaziga aylanishi zarur”, 17 July 2026. \url{https://gov.uz/oz/news/view/194551}.
\bibitem{chile} Ministerio de Desarrollo Social y Familia (Chile), División de Focalización. Orientaciones para el cálculo de la Calificación Socioeconómica. June 2024, pp. 4, 18. \url{https://www.registrosocial.gob.cl/docs/Orientaciones_para_el_calculo_de_la_Calificacion_Socieconomica.pdf}.
\bibitem{wb2025} Guven M., Yeachuri A., Almenfi M. Global Insights on Social Registries: Coverage and Beyond. World Bank, June 2025, pp. 3, 13, 25–26.
\bibitem{portal} Portal for discussion of draft normative legal acts. List of drafts placed for discussion from 11 September to 1 October 2026; checked 1 October 2026. \url{https://regulation.adliya.uz}.
\end{thebibliography}\endgroup
\end{document}
